%%
%% This is file `acmlarge.tex',
%% generated with the docstrip utility.
%%
%% The original source files were:
%%
%% samples.dtx  (with options: `all,acmlarge')
%% 
%% IMPORTANT NOTICE:
%% 
%% For the copyright see the source file.
%% 
%% Any modified versions of this file must be renamed
%% with new filenames distinct from acmlarge.tex.
%% 
%% For distribution of the original source see the terms
%% for copying and modification in the file samples.dtx.
%% 
%% This generated file may be distributed as long as the
%% original source files, as listed above, are part of the
%% same distribution. (The sources need not necessarily be
%% in the same archive or directory.)
%%
%%
%% Commands for TeXCount
%TC:macro \cite [option:text,text]
%TC:macro \citep [option:text,text]
%TC:macro \citet [option:text,text]
%TC:envir table 0 1
%TC:envir table* 0 1
%TC:envir tabular [ignore] word
%TC:envir displaymath 0 word
%TC:envir math 0 word
%TC:envir comment 0 0
%%
%% techbrief-sample.tex -- an ACM TechBrief, typeset with the unmodified
%% acmart class plus techbrief.sty.  Derived from the acmart sample
%% acmlarge.tex (single column, US Letter); the content reproduces
%% ACM TechBrief Issue 18 (Spring 2026), doi:10.1145/3819832.
%%
%% Build with `make` (LuaLaTeX + BibTeX via latexmk).
%%
%% The `nonacm' option switches off the ACM journal furniture (reference
%% format block, CCS concepts, copyright box); techbrief.sty supplies
%% the TechBrief banner, running heads, frame and back page instead.
%% `natbib=false' lets techbrief.sty use acmart's biblatex style
%% (acmnumeric), with references numbered in citation order and \cite
%% producing superscripts.
\documentclass[acmlarge,nonacm,natbib=false]{acmart}
\usepackage{techbrief}

%% Bibliography sources (one \addbibresource command per source).
\addbibresource{techbrief-sample.bib}

%% Bibliographic information, printed on the back page.
\acmDOI{10.1145/3819832}
\acmISBN{979-8-4007-2785-6}

%% Issue information, printed in the footer of every page.
\techbriefissue{Spring 2026}{18}

%% Acknowledgments, printed in the panel at the foot of the back page.
\techbriefacknowledgments{The authors would like to thank the following
  subject-matter experts for their reviews and input: Mr.~Gerhard
  Schimpf, Dr.~Benjamin Cedric Larsen, and Dr.~Geoff Cohen. A thank you
  to the ACM TechBrief Committee and the ACM Technology Policy Council
  for their valuable input. The PDF version of the document meets WCAG
  2.2 Level AA accessibility standards and complies with the PDF
  Accessibility Checker (PAC2024) requirements.}

%%
%% end of the preamble, start of the body of the document source.
\begin{document}

%%
%% The "title" command has an optional parameter, the short title used
%% in the running head of pages 2 onward.  Use \\ in the full title to
%% control where the banner title breaks.
\title[Autonomy, Opportunities, and Challenges]{Autonomy, Opportunities,
  and Challenges\\of Action-Taking AI Systems}

%%
%% Authors, in byline order.  Only the names are printed.
\author{Rohith Nama}
\author{Dr.~Larry Medsker}
\author{Brian Peretti}

\maketechbrief

\begin{techbriefbody}
Agentic AI systems include an AI component that plans and executes
sequences of actions toward a user-defined goal, accessing external
software, data, and services along the way. As capabilities continue to
advance, defining Agentic Systems will evolve. Earlier AI systems with a
chat interface produced text or other content for human review; agentic
systems can take actions with real-world consequences, such as browsing
the web, writing and executing code, sending messages to other agents or
humans, and managing files. The decision of how much human oversight and
control to exert over agentic AI systems depends on policy and design.
Even with oversight, users must delegate some authority to these systems
to act on their behalf.

%% \looseness=-1 is ordinary copyfitting: it asks TeX to set this
%% paragraph one line shorter, which keeps the first page's layout.
\looseness=-1
This shift from output-generating to action-taking systems introduces
risks and responsibilities that existing legal, regulatory, and technical
frameworks do not currently address. Adoption is accelerating: a 2025
survey of over 500 technology leaders found that 48\% are already
adopting or deploying agentic AI.\cite{ey2025pulse} With national and
sub-national laws coming online, agentic AI raises questions that these
frameworks, designed primarily for generative AI, do not fully address.

Agentic AI systems are typically built on large language models and
inherit the limitations of chat-based AI, including hallucination and
bias. Autonomy exists on a continuum: some require explicit authorization
for every action; others operate completely independently toward a
stated goal. Most require human approval for consequential actions, but
approval can be waived over time. The boundary between autonomous and
supervised actions is a design choice that can change without user
notification.

\begin{techbrieftable}{Differences between chat-based AI and agentic AI systems}
    {L{115.3bp}L{171.7bp}L{228.6bp}}
  \rowcolor{tbblue}
    & \tbhead{Chat-based AI} & \tbhead{Agentic AI} \tbheadend
  \tbrowhead{Primary mechanism}
    & Produces content in response to a prompt
    & Plans and executes multi-step tasks to achieve a goal \\
  \tbrowhead{Human role}
    & User reviews output before acting
    & User sets goal; agent acts without step-by-step approval \\
  \tbrowhead{External access}
    & Limited; may browse the web, but does not write to external systems
    & Reads from and writes to external systems with software interfaces \\
  \tbrowhead{Error recovery}
    & User has the opportunity to review the output before acting on it
    & Agents may act on an error before human review. Agents may provide
      an ``undo'' feature that may attempt to revert external changes. \\
  \tbrowhead{Accountability}
    & Typically attributed to the user who acts on the output
    & Potentially distributed across model provider, framework, deployer,
      and user \\
  \tbrowhead{Primary security concern}
    & Hallucination, bias, insecurely generated code, and reproduction of
      copyrighted content
    & Prompt injection, tool misuse, cascading failures, plus concerns
      from Chat-based systems. \\
\end{techbrieftable}

\section*{Key Policy Dimensions}

\runinhead{Legal Responsibility and Liability.} Agentic AI users
typically define high-level goals without specifying operational
boundaries, such as stop conditions or escalation triggers.
Accountability for harms caused by these systems is thus diffused in ways
that existing legal doctrines were not designed to handle. Responsibility
may be attributed to the model provider, framework developer, deployer,
or end user, as each shapes the system's behavior, but none controls any
specific action. This makes it difficult to assess responsibility for
harms after the fact.

\begin{itemize}
\item In one incident, an AI coding agent deleted a company's entire
  production database despite explicit instructions to make no changes.
  Responsibility for the losses might be attributed to the company that
  deployed the agent, the platform provider, the model provider, or the
  developer who commanded the agent and configured its permissions. No
  person made ``the decision'' to delete the database, yet harm was
  done.\cite{nolan2025replit}
\item There is no case law applying existing doctrines such as product
  liability, negligence, and software licensing to Agentic AI. Software
  licensing allocates risk through contractual terms, but these may not
  map to systems that take autonomous actions affecting third parties not
  party to the license.\cite{bertolini2025liability}
\item The European Parliament's 2025 study on AI civil liability found
  that existing product liability rules struggle with autonomous AI
  systems, making causation and fault difficult to establish across AI
  supply chains.
\item Accountability may also depend on uneven infrastructure across
  deployments, including audit trails, replayable action logs, and
  standardized incident reporting. Where such infrastructure is absent,
  establishing what happened and why is difficult, regardless of how
  liability rules are framed.
\end{itemize}

\begin{takeaway}
Resolving liability questions for agentic AI will likely require both
substantive legal clarification and parallel investment in the technical
mechanisms needed to apply it.
\end{takeaway}

\runinhead{Security and Operational Risk.} Agentic AI introduces
security challenges beyond those of chat-based systems.

\begin{itemize}
\item Because LLMs process text as both data and commands, Agentic AI
  systems cannot reliably distinguish between legitimate content and
  embedded instructions.\cite{perez2022ignore}
  \begin{itemize}
  \item In one case, malicious instructions embedded in a Slack message
    caused an AI agent to retrieve and expose content from private
    channels without any user action.\cite{promptarmor2024slack}
  \end{itemize}
\item Agents that retain memory across sessions may also be vulnerable to
  memory poisoning, where malicious inputs persist and influence future
  actions.
\item Agents are often configured with broad access rather than
  permissions scoped to specific actions, which can increase exposure
  when errors occur.
\item Additional concerns include vulnerabilities in agent-generated
  code, unauthorized access to copyrighted data, and risks from
  third-party extensions and plugins.
\item Although guardrail frameworks that monitor agent behavior and
  detect prompt injection are under development, none currently provide
  comprehensive protection.\cite{rebedea2023nemo,chennabasappa2025llamafirewall}
\end{itemize}

\begin{takeaway}[12bp]
Agentic systems' security depends on choices about permissions,
oversight, and the boundaries of autonomous action, choices that are
both technical and organizational.
\end{takeaway}

\runinhead{Consumer Transparency and Recourse.} When an AI agent acts on
an individual's behalf, the user may not know what systems the agent can
connect to, what actions it can take without explicit confirmation, or
how to revoke permissions once granted. No standardized disclosure format
currently exists for agent permissions.

\begin{itemize}
\item OpenClaw can read and send emails, manage calendars, and execute
  terminal commands through consumer messaging apps, accumulating more
  than 300,000 users within weeks of its 2025
  release.\cite{zaitsev2026openclaw} Security researchers found that
  approximately 12\% of third-party extensions in its marketplace
  contained malicious code, including keyloggers and data exfiltration
  tools.\cite{chang2026openclaw}
\item Because many agents communicate in natural language and exhibit
  human-like responsiveness, users may develop trust in their
  reliability, particularly in sensitive domains such as healthcare,
  where an agent's confident tone can mask clinical
  uncertainty.\cite{bellogin2025systemic}
\item Agents that act through natural language interfaces can also
  introduce language-mediated harms. An agent that sends an email or
  schedules a meeting on a user's behalf may damage relationships,
  reputation, or trigger legal exposure.
\end{itemize}

\begin{takeaway}
Consumer transparency for agentic AI involves disclosing that an AI is in
use, making the scope of an agent's authority visible, and accounting for
potential harms.
\end{takeaway}

\runinhead{Multi-Agent Coordination Risks.} When multiple agents interact
in coordinated networks, they can produce unexpected outcomes.
Multi-agent networks may also reinforce one another's errors, a pattern
called conformity bias.\cite{schroederdewitt2025multiagent,ashery2025emergent}

\runinhead{Workforce and Societal Impact.} Organizations view agentic AI
as a means to reduce current headcount, particularly in entry- and
mid-level roles. A 2025 Gartner survey of supply chain leaders found that
55\% expect agentic AI to reduce hiring for entry-level positions, and
51\% believe the technology will drive a shift to overall workforce
reductions.\cite{gartner2026survey} Similar dynamics are happening across
multiple sectors.

\begin{itemize}
\item Reductions in entry-level positions will likely shrink pathways
  through which workers traditionally develop expertise. Where junior
  roles disappear, the path to mid-career and senior expertise narrows,
  raising questions about how skills and judgment are formed in fields
  where hands-on practice has been essential.
\item Independent studies measuring the actual effects of agentic AI on
  organizational effectiveness, service quality, employment composition,
  or equity remain limited. Productivity claims have not been
  independently verified at scale, even as those claims are already
  shaping hiring decisions.
\end{itemize}

\begin{takeaway}[12bp]
Accelerating displacement, narrowing entry pathways, and unverified
productivity gains may have long-term effects on skill formation and
labor markets that are not yet measurable but are already shaping
organizational behavior. Tracking these effects systematically may be a
precondition for informed labor and education policy responses.
\end{takeaway}

\section*{Conclusion}

Agentic AI represents a shift from AI systems that generate content for
human review to systems that take actions on humans' behalf. This raises
questions about legal liability, consumer transparency, and regulatory
adequacy that are distinct from earlier generations of AI.

Many technical and operational challenges are being addressed in
industries through risk management, engineering practices, and emerging
standards. Relevant standards efforts are underway, including NIST's AI
Agent Standards Initiative, OECD analytical frameworks, and
industry-funded research into secure agent ecosystems.\cite{nist2026caisi}
Yet while enterprise deployments often involve staged integration and
constrained environments, consumer-facing systems are progressing with
less structured oversight.

In sectors such as healthcare, financial services, and critical
infrastructure, where existing frameworks assume a human decision-maker,
existing law may need to be supplemented by sector-specific guidance
tailored to autonomous systems. Agentic systems that operate across
jurisdictions for a single task also raise cross-border accountability
questions that existing international mechanisms have yet to address.

A central policy question is who is responsible when autonomous systems
cause harm, and how to ensure adequate consumer transparency when agents
act on an individual's behalf. The technical question is how to make AI
Agentic Systems reliably follow the policies that they are given. To
align these goals, safeguards such as defined authentication and
delegation, the ability to audit and monitor systems, support for
investigations and fault determination, and further research into
inter-agent dynamics are needed.\cite{schmidt2025secure}
\end{techbriefbody}

%%
%% The back page: references, with the "Additional Information" and
%% acknowledgments panel at the foot.
\techbriefbackpage
\printbibliography

\end{document}
\endinput
%%
%% End of file `techbrief-sample.tex'.
